% ============================================================================
% CAPÍTULO 6: COSMOLOGÍA
% ============================================================================
\chapter{Cosmología}
\label{cap:cosmologia}

\begin{center}
\textit{``El universo a gran escala''}
\end{center}

\vspace{0.5cm}

\begin{tcolorbox}[colback=white, colframe=kleinblue, boxrule=2pt]
\centering
\textit{``El universo no solo tiene una historia,\\
tiene una TEMPERATURA.''}\\
--- Cosmología moderna\\[0.3cm]
\textit{``Y esa temperatura está escrita en $7\pi$.''}\\
--- Teoría Klein, 2026
\end{tcolorbox}


% ============================================================================
\section{El Universo a Gran Escala}
% ============================================================================

El universo comenzó hace 13.8 mil millones de años en el Big Bang.
Desde entonces:
\begin{itemize}
    \item Se expandió (y sigue expandiéndose)
    \item Se enfrió (de $10^{32}$ K a 2.7 K)
    \item Formó estructuras (galaxias, estrellas, planetas)
\end{itemize}

Tres observaciones clave conectan con Klein:
\begin{enumerate}
    \item La temperatura del CMB (radiación cósmica de fondo)
    \item La constante cosmológica (energía oscura)
    \item El número de Avogadro (conexión macro-micro)
\end{enumerate}


% ============================================================================
\section{La Constante Cosmológica: $\rho_\Lambda/\rho_P = (7/2) \times (7\pi)^{-92}$}
% ============================================================================

(Ya visto en Capítulo 4, aquí profundizamos)

\begin{tcolorbox}[colback=kleingray, colframe=kleinblue]
El problema de la constante cosmológica es el ``peor desacuerdo''
en la historia de la física: $10^{123}$ órdenes de magnitud.

\textbf{FÓRMULA KLEIN:}

\begin{equation}
\boxed{\frac{\rho_\Lambda}{\rho_P} = \frac{7}{2} \times (7\pi)^{-92}}
\end{equation}

donde $92 = 4 \times 23 = 4D \times (\dim(\text{SU}(5)) - 1)$

$(7\pi)^{-92} \approx 10^{-124}$

¡Esto EXPLICA los 123 órdenes de magnitud!\\
Error: $\sim$0.7\%
\end{tcolorbox}


% ============================================================================
\section{La Temperatura del CMB: $T_{\text{CMB}} = \pi T_P / (7\pi)^{24}$}
% ============================================================================

\begin{nivelconceptual}
\textbf{La Radiación Cósmica de Fondo (CMB)}

380,000 años después del Big Bang, el universo se enfrió lo suficiente
para que los electrones se unieran a los protones formando hidrógeno.

La luz de ese momento TODAVÍA llena el universo.
La vemos como radiación de microondas a 2.7255 K.

Esta es la ``foto más antigua'' del universo.
\end{nivelconceptual}

\begin{nivelmatematico}
\textbf{$T_{\text{CMB}}$ desde Klein}

\begin{equation}
\boxed{T_{\text{CMB}} = \frac{\pi \times T_{\text{Planck}}}{(7\pi)^{24}}}
\end{equation}

donde:
\begin{itemize}
    \item $T_{\text{Planck}} = 1.417 \times 10^{32}$ K (temperatura de Planck)
    \item $24 = \dim(\text{SU}(5)) = 5^2 - 1$
\end{itemize}

\textbf{VERIFICACIÓN:}
\begin{align}
\text{Predicción:} & \quad T_{\text{CMB}} = 2.7315 \text{ K}\\
\text{Observado:} & \quad T_{\text{CMB}} = 2.7255 \text{ K (Planck 2018)}\\
\text{Error:} & \quad 0.22\%
\end{align}

\textbf{INTERPRETACIÓN:}

El universo se ha enfriado por un factor de $(7\pi)^{24}$ desde el Big Bang.
El exponente $24 = \dim(\text{SU}(5))$ conecta cosmología con física de partículas.
\end{nivelmatematico}


% ============================================================================
\section{El Número de Avogadro: $N_A = e^{(5/2 - 1/99) \times 7\pi}$}
% ============================================================================

\begin{nivelconceptual}
\textbf{El Puente Macro-Micro}

$N_A = 6.022 \times 10^{23}$ es el número de átomos en un mol de sustancia.

Conecta el mundo macroscópico (gramos, litros) con el microscópico
(átomos, moléculas).

¿Por qué este número específico?
\end{nivelconceptual}

\begin{nivelmatematico}
\textbf{$N_A$ desde Klein}

\begin{equation}
\boxed{N_A = e^{(5/2 - 1/99) \times 7\pi}}
\end{equation}

\textbf{Desglose:}
\begin{itemize}
    \item $5/2 = 2.5$ (corrección termodinámica, 5 dimensiones)
    \item $1/99 \approx 0.0101$ (corrección fina)
    \item $5/2 - 1/99 = 2.4899$
    \item $\times 7\pi = \times 21.991$
    \item Exponente total $= 54.76$
\end{itemize}

\textbf{VERIFICACIÓN:}
\begin{align}
\text{Predicción:} & \quad N_A = 6.027 \times 10^{23}\\
\text{Observado:} & \quad N_A = 6.022 \times 10^{23}\\
\text{Error:} & \quad 0.08\%
\end{align}
\end{nivelmatematico}


% ============================================================================
\section{Resumen del Capítulo}
% ============================================================================

\begin{tcolorbox}[colback=kleingray, colframe=kleinblue, title=\textbf{Resumen del Capítulo 6}]

\textbf{PREDICCIONES COSMOLÓGICAS KLEIN:}

\begin{enumerate}
    \item $\rho_\Lambda/\rho_P = (7/2) \times (7\pi)^{-92}$ \hfill [0.7\% error] --- Constante cosmológica
    \item $T_{\text{CMB}} = \pi T_P / (7\pi)^{24}$ \hfill [0.22\% error] --- Temperatura CMB
    \item $N_A = e^{(5/2-1/99) \times 7\pi}$ \hfill [0.08\% error] --- Número de Avogadro
    \item $t_U/t_P = (7\pi)^{45}$ \hfill [$\sim$5\% error] --- Edad del universo
\end{enumerate}

\vspace{0.3cm}
\textbf{EXPONENTES CLAVE:}
\begin{itemize}
    \item $24 = \dim(\text{SU}(5)) \to T_{\text{CMB}}$
    \item $45 \approx 2 \times 24 - 3 \to$ edad del universo
    \item $92 = 4 \times 23 \to$ constante cosmológica
\end{itemize}

\vspace{0.3cm}
\textbf{PRÓXIMO CAPÍTULO:} Antimateria

\end{tcolorbox}


% ============================================================================
\section*{Código de Verificación}
% ============================================================================

\begin{lstlisting}[caption={Verificación numérica - Capítulo 6}]
import numpy as np
from numpy import pi

siete_pi = 7 * pi

# Temperatura del CMB
T_P = 1.416808e32  # K
T_CMB_obs = 2.7255  # K
T_CMB_klein = pi * T_P / siete_pi**24
print(f"TEMPERATURA CMB:")
print(f"  T_CMB Klein = {T_CMB_klein:.4f} K")
print(f"  T_CMB obs   = {T_CMB_obs} K")
print(f"  Error = {abs(T_CMB_klein - T_CMB_obs)/T_CMB_obs*100:.2f}%")

# Numero de Avogadro
N_A_obs = 6.02214076e23
coef = 5/2 - 1/99
N_A_klein = np.exp(coef * siete_pi)
print(f"\nNUMERO DE AVOGADRO:")
print(f"  N_A Klein = {N_A_klein:.5e}")
print(f"  N_A obs   = {N_A_obs:.5e}")
print(f"  Error = {abs(N_A_klein - N_A_obs)/N_A_obs*100:.2f}%")
\end{lstlisting}
